\chapter{Modelling Overnight-Rate Products}\label{ch:rc:modelling-overnight-rate-products}

A company's floating-rate loan pays, each quarter, SOFR compounded daily over
the quarter, known only when the quarter ends. Its treasurer buys a cap, and
the desk that sells it notices something the old interbank caps never did: the
first caplet, whose period has already started, still has optionality, and
every caplet keeps accruing risk after its period begins, until its last
overnight fixing. An interbank rate fixed at the start of its period; a
compounded overnight rate fixes at the end. The models of chapters~7 and~8 were
built for rates known in advance. This chapter extends them to rates known in
arrears: in the \omterm{def:rc:short-rate-models:hw}{Hull--White model} the \omterm{def:rc:modelling-overnight-rate-products:caplet}{backward-looking caplet} has a closed
form, in the market model the extension is the \omterm{def:rc:modelling-overnight-rate-products:gfmm}{generalised forward market model},
and with the calendar of central-bank meetings the rate becomes a step function
whose risk arrives on meeting days.

\section{Backward-looking rates and their options}

\begin{definition}[Backward-looking rate, forward-looking term rate]\label{def:rc:modelling-overnight-rate-products:rates}
A \emph{backward-looking rate}\index{backward-looking rate} for a period $[S,E]$ is
the overnight rate compounded in arrears over the period (One Quant Book 2,
chapter 1): $1+\delta R(S,E) = \prod_i(1+r_id_i/360)$, known only at $E$. A
\emph{forward-looking term rate}\index{forward-looking term rate} for the same
period is published at $S$, as an expectation of that compounded rate implied by
futures and swaps (a term SOFR rate) or as an interbank offered rate (Euribor);
it is known at the start.
\end{definition}

\begin{definition}[Backward-looking caplet]\label{def:rc:modelling-overnight-rate-products:caplet}
A \emph{backward-looking caplet}\index{backward-looking caplet} pays
$\delta\max(R(S,E)-K,0)$ at $E$ (or a few days later) on the compounded rate of
its own period.
\end{definition}

\begin{omfigure}
\begin{tikzpicture}[x=2.0cm,font=\small]
\draw[->,thick] (0,0) -- (5.4,0) node[right] {time};
\foreach \x/\l in {0/0,2/$S$,4/$E$} {\draw (\x,-0.1) -- (\x,0.1); \node[below] at (\x,-0.12) {\l};}
\draw[omDef,thick,fill=omDef!12] (2,0.4) rectangle (4,0.8); \node[font=\footnotesize] at (3,0.6) {accrual period};
\draw[omThm,very thick] (0,1.4) -- (2,1.4); \draw[omThm,thick,fill=omThm] (2,1.4) circle (2pt);
\node[text=omThm,anchor=west,font=\footnotesize] at (4.15,1.4) {forward-looking: risk until $S$, rate known at $S$};
\draw[omProp,very thick] (0,2.1) -- (4,2.1); \draw[omProp,thick,fill=omProp] (4,2.1) circle (2pt);
\node[text=omProp,anchor=west,font=\footnotesize] at (4.15,2.1) {backward-looking: risk until $E$, rate known at $E$};
\end{tikzpicture}

\omcaption{A forward-looking rate carries risk until the start of its period
and is known from then on; a compounded overnight rate keeps moving until the
last fixing at the end, so an option on it carries risk through its whole
accrual period.}\label{fig:rc:modelling-overnight-rate-products:timing}
\end{omfigure}

\begin{proposition}[Backward-looking caplets in Hull--White]\label{prop:rc:modelling-overnight-rate-products:hw}
Approximate daily compounding by $1+\delta R = \exp\int_S^Er_u\,du$. In Hull--White
with constant $\sigma$ and $c = 1+\delta K$,
\[
\mathrm{Caplet}_0 = P(0,S)\bigl[\Phi(-d_2) - c\,m\,\Phi(-d_1)\bigr],\quad
m = \frac{P(0,E)}{P(0,S)},\ d_1 = \frac{\ln(mc)+V/2}{\sqrt V},\ d_2 = d_1-\sqrt V,
\]
with $V = V_{\text{before}}+V_{\text{during}}$,
$V_{\text{before}} = \sigma^2B(\delta)^2\frac{1-e^{-2\kappa S}}{2\kappa}$ and
$V_{\text{during}} = \sigma^2\int_0^\delta B(\tau)^2\,d\tau\approx\sigma^2\delta^3/3$. The
forward-looking caplet on a term rate for the same period has the same formula
with $V = V_{\text{before}}$ only.
\end{proposition}
\begin{proof}
The payoff is $(e^{-\int_0^Sr}-c\,e^{-\int_0^Er})^+$. Under the measure with density
$e^{-\int_0^Sr}/P(0,S)$ the price is $P(0,S)\E[(1-cL)^+]$ with
$L=e^{-\int_S^Er}$, lognormal with mean $P(0,E)/P(0,S)$; its log-variance is the
variance of $\int_S^Ex_u\,du$, which splits into the shocks before $S$ (through the
bond price at $S$) and those during the period. The expectation is a Black put.
\end{proof}

\omcode{code/firm/rfrcaplet/firm_rfrcaplet.py}{35}{56}{The two parts of the variance, and the backward-looking caplet in closed form.}\label{lst:rc:modelling-overnight-rate-products:hw}

\begin{example}[A two-year SOFR cap]\label{ex:rc:modelling-overnight-rate-products:cap}
On chapter~1's SOFR curve with $\kappa=3\%$ and $\sigma=90$ basis points, a two-year cap on
three-month compounded SOFR at 3.75\% on USD~100 million costs USD~288\,963: eight
caplets from USD~13\,623 (the current quarter) to 53\,846
(\cref{fig:rc:modelling-overnight-rate-products:caplets}). The same cap on a
three-month term rate, whose first caplet has already fixed, costs USD~250\,741:
the backward-looking cap is USD~38\,222 more expensive, of which 13\,623 is the
first caplet and 24\,599 the in-period variance of the other seven. A Monte Carlo
of the short rate with 260 steps a year reprices the caplet on the fifth quarter
(USD~40\,944) to within a fifth of a standard error.
\end{example}

\begin{omfigure}
\begin{tikzpicture}
\begin{axis}[width=12cm,height=5.4cm,ybar,bar width=10pt,
  ylabel={caplet value (USD thousand)},xlabel={period (years from spot)},
  tick label style={font=\small},label style={font=\small},
  xmin=-0.6,xmax=7.6,ymin=0,ymax=60,grid=major,grid style={gray!25},
  xtick=data,xticklabels from table={figdata/rates-credit-risk/10-modelling-overnight-rate-products/caplets.csv}{label},
  x tick label style={font=\scriptsize},
  legend columns=2,legend style={font=\footnotesize,at={(0.5,-0.3)},anchor=north,draw=none,fill=none,/tikz/every even column/.append style={column sep=8pt}},
  table/col sep=comma]
\addplot[fill=omProp!55,draw=omProp] table[x=i,y=back]{figdata/rates-credit-risk/10-modelling-overnight-rate-products/caplets.csv};
\addplot[fill=omDef!55,draw=omDef] table[x=i,y=fwd]{figdata/rates-credit-risk/10-modelling-overnight-rate-products/caplets.csv};
\legend{backward-looking (compounded SOFR),forward-looking (term rate)}
\end{axis}
\end{tikzpicture}

\omcaption{Caplets of a two-year cap at 3.75\% on USD~100 million, on compounded
SOFR and on a term rate for the same quarters. The current quarter has a caplet
only in the backward-looking cap; each later \omterm{def:rc:modelling-overnight-rate-products:caplet}{backward-looking caplet} is worth more
by its period's in-arrears variance. Data: chapter~1's curve, Hull--White $\kappa=3\%$,
$\sigma=90$ basis points; the chapter's tutorial.}\label{fig:rc:modelling-overnight-rate-products:caplets}
\end{omfigure}

The in-period share of the variance is 100\% for the current quarter, 25\% for
the next, 7.9\% for the fifth and 4.8\% for the eighth: close to $(\delta/3)/(S+\delta/3)$,
the ratio a linear decay of volatility through the period gives
(\cref{fig:rc:modelling-overnight-rate-products:share}).

\begin{omfigure}
\begin{tikzpicture}
\begin{axis}[width=11cm,height=5cm,
  xlabel={start of the quarter $S$ (years)},ylabel={in-period share of variance (\%)},
  tick label style={font=\small},label style={font=\small},
  xmin=0,xmax=1.8,ymin=0,ymax=105,grid=major,grid style={gray!25},
  legend columns=2,legend style={font=\footnotesize,at={(0.5,-0.25)},anchor=north,draw=none,fill=none,/tikz/every even column/.append style={column sep=8pt}},
  table/col sep=comma]
\addplot[omThm,thick,mark=*] table[x=s,y=hw]{figdata/rates-credit-risk/10-modelling-overnight-rate-products/share.csv};
\addplot[omDef,thick,dashed] table[x=s,y=lin]{figdata/rates-credit-risk/10-modelling-overnight-rate-products/share.csv};
\legend{Hull--White,linear decay: $(\delta/3)/(S+\delta/3)$}
\end{axis}
\end{tikzpicture}

\omcaption{Share of a \omterm{def:rc:modelling-overnight-rate-products:caplet}{backward-looking caplet}'s variance that accrues during its own
quarter, by the quarter's start: all of it for the current quarter, a quarter for the
next, a few per cent two years out. Data: the chapter's
tutorial.}\label{fig:rc:modelling-overnight-rate-products:share}
\end{omfigure}

\section{The generalised forward market model}

\begin{definition}[Generalised forward market model]\label{def:rc:modelling-overnight-rate-products:gfmm}
The \emph{generalised forward market model}\index{generalised forward market model}
(Lyashenko and Mercurio, 2019) models, for each accrual period $[T_{k-1},T_k]$,
the forward $R_k(t) = \E^{T_k}_t[R(T_{k-1},T_k)]$ of the \omterm{def:rc:modelling-overnight-rate-products:rates}{backward-looking rate},
for all $t\le T_k$: before $T_{k-1}$ it coincides with the forward-looking forward
of chapter~8, and during the period it keeps moving with the overnight fixings. Its
volatility is $\sigma_k(t)g_k(t)$ with $g_k=1$ before $T_{k-1}$ and $g_k$ decaying to
zero at $T_k$, for instance linearly.
\end{definition}

\begin{proposition}[Effective variance with linear decay]\label{prop:rc:modelling-overnight-rate-products:gfmm}
With constant $\sigma_k$ and $g_k(t) = (T_k-t)/(T_k-T_{k-1})$ on the period, the
\omterm{def:rc:modelling-overnight-rate-products:caplet}{backward-looking caplet} is priced by Black's formula with total variance
$\sigma_k^2\bigl(T_{k-1}+\tfrac13(T_k-T_{k-1})\bigr)$, against $\sigma_k^2T_{k-1}$ for the
forward-looking caplet.
\end{proposition}
\begin{proof}
$\int_0^{T_{k-1}}\sigma_k^2\,dt+\int_{T_{k-1}}^{T_k}\sigma_k^2\bigl(\frac{T_k-t}{\delta}\bigr)^2dt =
\sigma_k^2(T_{k-1}+\delta/3)$, and $R_k$ is lognormal under $\mathbb Q^{T_k}$.
\end{proof}

The generalised model contains the old one: one process gives both the
forward-looking and the \omterm{def:rc:modelling-overnight-rate-products:rates}{backward-looking rates}, so a book of legacy interbank
products and new overnight ones is risk-managed in one model, with the market
model's simulation and calibration machinery.

\section{Low-dimensional Markov models}

Market models are high-dimensional; production systems for exotics and
exposures prefer models with a handful of state variables whose forward curve
has an exact formula.

\begin{definition}[Cheyette model]\label{def:rc:modelling-overnight-rate-products:cheyette}
A \emph{Cheyette model}\index{Cheyette model} (1992) is an HJM model whose forward
volatility factorises as $\sigma_f(t,T) = \sigma_r(t,x_t,y_t)\,e^{-\kappa(T-t)}$. Then
$f(t,T) = f(0,T)+e^{-\kappa(T-t)}\bigl(x_t+B(t,T)y_t\bigr)$ with two state variables,
$dx = (y-\kappa x)\,dt+\sigma_r\,dW$ and $dy = (\sigma_r^2-2\kappa y)\,dt$. With
deterministic $\sigma_r$ it is Hull--White; letting $\sigma_r$ depend on $x$ gives a
local-volatility (skewed) \omterm{def:rc:short-rate-models:short}{short-rate model} that is still Markov in $(x,y)$.
\end{definition}

The integral of the short rate over an accrual period is one more state variable
in such a model, so backward-looking products stay low-dimensional: this is the
Markov route to the same caplets, and the one that PDE and tree pricers follow.

\section{Meeting dates and jumps in the short rate}

The overnight rate does not diffuse: it sits at the policy rate and moves on the
days a central bank decides (One Quant Book 2, chapter 8). Its uncertainty
arrives in lumps, on a published calendar.

\begin{dated}{2026-09}{Meeting calendar and term rates}\label{dat:rc:modelling-overnight-rate-products:term}
The Federal Reserve's policy committee meets on 27--28 October and 8--9 December
2026, and on 26--27 January, 16--17 March, 27--28 April, 8--9 June, 27--28 July,
14--15 September, 26--27 October and 7--8 December 2027 (tentative until confirmed
at the preceding meeting). The ARRC formally recommended CME Group's forward-looking
term SOFR rates on 29 July 2021, and its best practices do not support their use in
derivatives except by end users hedging cash products that reference them.
\end{dated}

\omcode{code/firm/rfrcaplet/firm_rfrcaplet.py}{106}{111}{The variance of a period's average rate when the rate moves only at meetings.}\label{lst:rc:modelling-overnight-rate-products:meetings}

\begin{example}[One quarter, two meetings]\label{ex:rc:modelling-overnight-rate-products:meetings}
The quarter from 29 September to 29 December 2027 (years 1.0 to 1.25 from spot)
contains the meetings of 27 October and 8 December. If the rate moves only at
meetings, a decision on 27 October moves the quarter's average rate by 69\% of the
jump (the share of the quarter left), one on 8 December by 23\%; after 8 December the
quarter's rate is known, three weeks before its end. The generalised model's linear
decay spreads the same risk smoothly
(\cref{fig:rc:modelling-overnight-rate-products:profile}): scaled to the same total,
the two agree before the quarter and differ within it.
\end{example}

\begin{omfigure}
\begin{tikzpicture}
\begin{axis}[width=12cm,height=5.4cm,
  xlabel={time $t$ (years from spot)},ylabel={remaining s.d. (\% of total)},
  tick label style={font=\small},label style={font=\small},
  xmin=0,xmax=1.25,ymin=0,ymax=105,grid=major,grid style={gray!25},
  legend columns=2,legend style={font=\footnotesize,at={(0.5,-0.25)},anchor=north,draw=none,fill=none,/tikz/every even column/.append style={column sep=8pt}},
  table/col sep=comma]
\addplot[omDef,very thick] table[x=t,y=gfmm]{figdata/rates-credit-risk/10-modelling-overnight-rate-products/profile.csv};
\addplot[omThm,thick,const plot] table[x=t,y=meet]{figdata/rates-credit-risk/10-modelling-overnight-rate-products/profile.csv};
\legend{generalised forward market model (linear decay),rate moves only at meetings}
\draw[gray,dashed] (axis cs:1.0,0) -- (axis cs:1.0,105) node[pos=0.95,left,font=\footnotesize,text=black] {period starts};
\end{axis}
\end{tikzpicture}

\omcaption{Remaining standard deviation, seen at $t$, of the compounded rate of the
quarter from year 1.0 to 1.25, in two models scaled to the same total: a smooth
decay through the period, or steps at the Federal Reserve's meeting days (ten
meetings before the quarter ends; the last inside it on 8 December 2027). Data:
Federal Reserve meeting calendar; the chapter's
tutorial.}\label{fig:rc:modelling-overnight-rate-products:profile}
\end{omfigure}

\begin{remark}[What the meeting view changes]\label{rem:rc:modelling-overnight-rate-products:meetings}
Prices change little, since the total variance is calibrated either way; risk
changes a lot. A caplet's theta and vega are concentrated on meeting days, and a
book of short-dated overnight options should be hedged meeting by meeting, with
the futures of One Quant Book 2, chapter 8, and options on them.
\end{remark}

\section{Forward-looking term rates}

Term rates survive for loans: a borrower wants to know the next quarter's
interest at the start. Derived from futures and overnight swaps, a term SOFR rate
is a published expectation; options and swaps on it are forward-looking
instruments of chapter~4's kind. The market keeps derivatives on the overnight rate
itself, where the liquidity is, and restricts term-rate derivatives to the hedges of
the loans that use them, so that the term rate does not come to rest on trading in
itself. The desk that sells a term-rate cap to a borrower therefore hedges it with
overnight-rate options and carries the basis between the two: the in-period
variance of \cref{ex:rc:modelling-overnight-rate-products:cap}.

\section{Tutorial: backward against forward}\label{tut:rc:modelling-overnight-rate-products:tutorial}

\begin{tutorial}
\textbf{Goal.} Price a two-year cap on compounded SOFR and on a term rate, check the
closed form by simulation, and compare smooth and meeting-date risk profiles.
\textbf{End state:}
\cref{fig:rc:modelling-overnight-rate-products:caplets,fig:rc:modelling-overnight-rate-products:profile}
and the numbers of \cref{ex:rc:modelling-overnight-rate-products:cap}.
\begin{tutsteps}
\item \textbf{Variances}: \texttt{hw\_period\_variances(kappa, sigma, S, E)} for each
quarter.
\item \textbf{Caps}: \texttt{cap\_table()} and \texttt{cap\_totals()}.
\item \textbf{Simulate}: \texttt{mc\_check()}, 260 steps a year so that every period
date is on the grid.
\item \textbf{Meetings}: \texttt{remaining\_sd\_profile()}; \texttt{fig\_rc\_rfr.py}
writes the charts.
\end{tutsteps}
\textbf{What to change next.} Run the simulation with 250 steps a year and an
unaligned grid (remove the check) and see the bias; price a caplet on a period
that contains no meeting.
\end{tutorial}

\section{Build: backward-looking caplets}\label{bld:rc:modelling-overnight-rate-products:rfrcaplet}

\begin{build}
\bfield{Purpose} Caps and floors on compounded overnight rates, the options of the
book's floating-rate loans since the end of the interbank rates.
\bfield{Interface} \texttt{hw\_period\_variances}; \texttt{hw\_backward\_caplet},
\texttt{hw\_forward\_caplet}, \texttt{hw\_backward\_caplet\_mc};
\texttt{gfmm\_effective\_variance}, \texttt{black\_caplet};
\texttt{meeting\_variance}, \texttt{bachelier\_caplet}.
\bfield{Rules} Continuous compounding of the short rate as the model of daily
compounding; simulation grids must contain every period date (enforced); numpy only.
\bfield{Acceptance tests} \texttt{code/firm/rfrcaplet/tests/}: a \omterm{def:rc:modelling-overnight-rate-products:caplet}{backward-looking
caplet} exceeds its forward-looking twin by its in-period variance, which is
$\sigma^2\delta^3/3$ for a period starting now; the closed form agrees with Monte Carlo
on an aligned grid and an unaligned grid is refused; the linear-decay variance and the
meeting variance have their closed forms.
\bfield{Stretch} Daily compounding with business-day calendars; lookback and payment
delay; a meeting-date \omterm{def:rc:short-rate-models:short}{short-rate model} calibrated to futures and their options.
\end{build}

\begin{omsources}
\item A.\ Lyashenko and F.\ Mercurio, ``Looking forward to backward-looking rates:
a modeling framework for term rates replacing LIBOR'', 2019.
\item O.\ Cheyette, ``Term structure dynamics and mortgage valuation'', \emph{Journal
of Fixed Income} 1(4), 1992.
\item ARRC, ``ARRC formally recommends term SOFR'', 29 July 2021.
\item Federal Reserve, \emph{Meeting calendars and information}.
\end{omsources}

\section{Exercises}

\begin{exercise}[$\star$]\label{exo:rc:modelling-overnight-rate-products:1}
Why does a cap on compounded SOFR include a caplet on the current quarter, when an
interbank cap never did?
\end{exercise}

\begin{exercise}[$\star$]\label{exo:rc:modelling-overnight-rate-products:2}
With linear decay, by how much is the total variance of a \omterm{def:rc:modelling-overnight-rate-products:caplet}{backward-looking caplet}
on $[2,2.25]$ larger than that of its forward-looking twin?
\end{exercise}

\begin{exercise}[$\star$]\label{exo:rc:modelling-overnight-rate-products:3}
In \cref{ex:rc:modelling-overnight-rate-products:meetings}, compute the weights
69\% and 23\% from the meeting dates.
\end{exercise}

\begin{exercise}[$\star\star$]\label{exo:rc:modelling-overnight-rate-products:4}
Using $V_{\text{during}}\approx\sigma^2\delta^3/3$ and $V_{\text{before}}\approx\sigma^2\delta^2S$,
estimate the in-period share of variance of the second quarter and compare with the
25\% of the text.
\end{exercise}

\begin{exercise}[$\star\star$]\label{exo:rc:modelling-overnight-rate-products:5}
Read \cref{fig:rc:modelling-overnight-rate-products:profile}. On which days does a
short caplet on the quarter lose most of its time value under the meeting view,
and what does that mean for its theta?
\end{exercise}

\begin{exercise}[$\star\star$]\label{exo:rc:modelling-overnight-rate-products:6}
Why do the best practices restrict derivatives on term SOFR, and what basis does a
desk that sells term-rate caps carry?
\end{exercise}

\begin{exercise}[$\star\star\star$]\label{exo:rc:modelling-overnight-rate-products:7}
\emph{Coding.} With \texttt{cap\_table()}, split the USD~38\,222 difference between
the two caps into the first caplet and the in-period variance of the others.
\end{exercise}

\begin{exercise}[$\star\star\star$]\label{exo:rc:modelling-overnight-rate-products:8}
\emph{Find the flaw.} ``We moved our cap pricer from LIBOR to SOFR by changing the
index name: it still skips the first caplet and uses the variance to each period's
start.''
\end{exercise}

\section{Problem: The SOFR Cap on a Leveraged Loan}

\begin{problem}[{Weekend problem --- a borrower chooses its cap}]\label{pb:rc:modelling-overnight-rate-products:1}
A private-equity-owned company has a USD~100 million two-year loan paying three-month
compounded SOFR plus a margin; its lenders require a cap at 3.75\%. Its bank offers a
cap on compounded SOFR and a cap on three-month term SOFR. Use chapter~1's curve and
Hull--White with $\kappa=3\%$ and $\sigma=90$ basis points.

\medskip\noindent\textbf{Part I --- The two caps.}
\begin{enumerate}
\item Give the price of each cap.
\item Why is the compounded-SOFR cap more expensive?
\item Give the value of its current-quarter caplet.
\item Which cap hedges the loan exactly, and what does the other leave?
\item Give the in-period share of variance for the second and the eighth quarters.
\end{enumerate}

\medskip\noindent\textbf{Part II --- Checks.}
\begin{enumerate}[resume]
\item Give the fifth quarter's caplet by the closed form and by simulation.
\item What went wrong when the simulation used 250 steps a year?
\item What does the \omterm{def:rc:modelling-overnight-rate-products:gfmm}{generalised forward market model}'s linear decay predict for the
fifth quarter's extra variance, and how does it compare with Hull--White's?
\item Why do both models give nearly the same answer?
\item Which inputs matter most for the price?
\end{enumerate}

\medskip\noindent\textbf{Part III --- Meetings.}
\begin{enumerate}[resume]
\item Which meetings fall in the fifth quarter?
\item When is its rate known under the meeting view?
\item What is the caplet's vega on the day after the December meeting?
\item How would the bank hedge the caplet meeting by meeting?
\item Does the meeting view change the price or the risk?
\end{enumerate}

\medskip\noindent\textbf{Part IV --- Judgement.}
\begin{enumerate}[resume]
\item Which cap would you advise the company to buy?
\item What basis risk does the bank take if it sells the term-rate cap?
\item Why do lenders accept either?
\item State the \emph{named result}: the two cap prices and the gap, with its split.
\item In one sentence: what does ``in arrears'' add to an option?
\end{enumerate}
\end{problem}

\section{Interview questions}

\begin{interviewq}[$\star$ \iqroles{trader, bank}]\label{iq:rc:modelling-overnight-rate-products:1}
What is the difference between a caplet on Euribor and a caplet on compounded
€STR?
\end{interviewq}

\begin{interviewq}[$\star\star$ \iqroles{researcher}]\label{iq:rc:modelling-overnight-rate-products:2}
How does the market model need to change for \omterm{def:rc:modelling-overnight-rate-products:rates}{backward-looking rates}?
\end{interviewq}

\begin{interviewq}[$\star\star$ \iqroles{researcher, developer}]\label{iq:rc:modelling-overnight-rate-products:3}
Price a \omterm{def:rc:modelling-overnight-rate-products:caplet}{backward-looking caplet} in Hull--White. Where does the extra variance come
from?
\end{interviewq}

\begin{interviewq}[$\star\star$ \iqroles{trader}]\label{iq:rc:modelling-overnight-rate-products:4}
How do central-bank meetings change the risk of short-dated overnight-rate options?
\end{interviewq}

\begin{interviewq}[$\star\star\star$ \iqroles{developer, risk}]\label{iq:rc:modelling-overnight-rate-products:5}
Your Monte Carlo price of a compounded-rate caplet is 6\% below the closed form with
a tight standard error. What do you check first?
\end{interviewq}

\begin{interviewq}[$\star\star\star$ \iqroles{researcher, bank}]\label{iq:rc:modelling-overnight-rate-products:6}
Why is term SOFR a problem as a derivatives underlying, and what do banks do instead?
\end{interviewq}
